
\subsubsection{6.1 Interpretation of Results}

Three feature types achieved 100\% test accuracy on the 4-model test set, exceeding the initial 70\% threshold. The perfect separation indicates strong discriminative signal for binary depth classification on this dataset. However, the small test set (n=4) provides limited statistical confidence. Under binomial assumptions, 4/4 correct yields a wide 95\% confidence interval approximately spanning 40\% to 100\%.

Random initialization tests showed zero performance degradation when using untrained models instead of pretrained models. This result indicates that the discriminative features reflect architectural properties (layer count, residual blocks, bottleneck ratios) that exist in the network structure regardless of training history. Features initially designed to capture gradient flow patterns (layer-wise norm slopes, depth-weighted aggregations) apparently functioned as proxies for architectural properties rather than measuring training-induced effects.

Within-family validation achieved 100\% accuracy for both ResNet and DenseNet families, indicating that depth signals persist even when controlling for architecture family differences. This suggests that the classification does not merely distinguish ResNet from VGG design patterns, but detects depth-related structural properties within architecture families.

\subsubsection{6.2 Feature Analysis}

The strongest discriminators identified through feature importance analysis were bottleneck ratio and layer count. Deep models (ResNet-50 and deeper, all DenseNet models) use bottleneck layers (1×1 → 3×3 → 1×1 convolution patterns) extensively, while shallow models (ResNet-18/34, all VGG variants) use standard 3×3 convolutions without bottlenecks. Layer count is directly related to the depth classification task by definition (shallow ≤34, deep ≥50).

Global weight statistics achieved perfect classification despite using only four aggregated values. Mean weight magnitude showed strong negative association with depth, with deep models exhibiting lower mean values (3-7 range) compared to shallow models (11-33 range). This pattern likely reflects parameter count scaling: more layers distribute parameters across more weight tensors, reducing per-tensor magnitude on average.

Batch normalization features achieved lower accuracy (75\%) compared to other feature types. Feature importance analysis suggests that batch normalization layer count (architectural property) drives classification rather than learned gamma/beta statistics. VGG models have zero batch normalization layers, creating a strong "0 BN → shallow" signal that may not generalize as robustly as comprehensive architectural features.

\subsubsection{6.3 Limitations}

\textbf{Sample Size.} The test set contains only 4 models, limiting statistical confidence in the 100\% accuracy estimate. Confidence intervals are wide, and larger-scale validation with 50-100 models would provide more precise accuracy estimates. Within-family validation provides some additional independent tests (ResNet n=9, DenseNet n=4), but each family subset also has small sample size.

\textbf{Task Scope.} The study tested binary classification (shallow vs deep) rather than continuous depth regression. Whether weight statistics enable predicting exact layer counts (18, 34, 50, 101, 152) or only coarse categorization remains untested. The 34-50 layer gap in the dataset avoids testing classification near the threshold boundary.

\textbf{Architecture Coverage.} The dataset includes only ResNet, VGG, and DenseNet families (2015-2017 era convolutional networks). Modern architectures remain untested:
\begin{itemize}
\item Vision Transformers (ViT, DEIT, Swin Transformer) use self-attention mechanisms rather than convolutions
\item Recent CNNs (EfficientNet, ConvNeXt, RegNet) incorporate compound scaling, inverted bottlenecks, or neural architecture search
\item Dynamically-routed networks with stochastic depth or adaptive computation may not exhibit fixed architectural signatures
\end{itemize}

Whether the findings generalize to these architectures is unknown.

\textbf{Confounding Variables.} Network depth, width, and architecture family are inherently correlated. Deep models are systematically wider (more channels per layer) and use different design patterns (bottleneck blocks, residual stages) compared to shallow models. Within-family validation partially controls for family-specific patterns, but complete isolation of depth effects independent of all other architectural properties is not feasible without synthetic architectures that vary only depth. Such synthetic architectures may not reflect real-world model properties, as very deep networks require skip connections for training stability.

\textbf{Dataset Generalization.} All models were pretrained on ImageNet-1K. Whether findings generalize to models pretrained on other datasets (CIFAR, COCO, domain-specific datasets) is untested. Random initialization tests suggest architectural features are training-dataset-independent, but empirical validation would strengthen this conclusion.

\textbf{Mechanism Interpretation.} While random initialization tests indicate features are architectural rather than training-induced, this does not exclude the possibility that alternative features designed to capture training effects might also enable classification. The tested features may not comprehensively span all possible training-induced patterns.

\subsubsection{6.4 Practical Implications}

Weight-based depth classification could enable architecture verification in model repositories without requiring metadata or forward passes. Given a model checkpoint file, weight statistics can be extracted and classified in seconds. Potential applications include deployment constraint validation, model repository quality control, and architecture claim verification.

Feature importance analysis identifies which architectural properties most strongly discriminate depth. Compression strategies that preserve these structural signatures (bottleneck layer structure, residual block organization) would maintain architectural detectability after compression.

The finding that architectural properties create detectable signatures independent of training suggests that other structural properties (network width, module counts, architecture family) might also be detectable from weight statistics alone, though this requires empirical validation.

\subsubsection{6.5 Comparison to Existing Methods}

Standard approaches to architecture detection include metadata parsing (requires complete and accurate metadata) and forward-pass-based methods (require test data and computational resources). The weight-based approach requires only checkpoint file access. Accuracy comparison to forward-pass methods is not available as prior work on depth classification from weight statistics was not identified.

The computational cost is substantially lower than forward-pass methods: feature extraction takes seconds with no GPU requirement, while forward passes scale with dataset size and require inference computation.
